% Options for packages loaded elsewhere
\PassOptionsToPackage{unicode}{hyperref}
\PassOptionsToPackage{hyphens}{url}
\PassOptionsToPackage{dvipsnames,svgnames,x11names}{xcolor}
%
\documentclass[
  10pt,
  a4paper,
]{article}
\usepackage{amsmath,amssymb}
\usepackage{iftex}
\ifPDFTeX
  \usepackage[T1]{fontenc}
  \usepackage[utf8]{inputenc}
  \usepackage{textcomp} % provide euro and other symbols
\else % if luatex or xetex
  \usepackage{unicode-math} % this also loads fontspec
  \defaultfontfeatures{Scale=MatchLowercase}
  \defaultfontfeatures[\rmfamily]{Ligatures=TeX,Scale=1}
\fi
\usepackage{lmodern}
\ifPDFTeX\else
  % xetex/luatex font selection
\fi
% Use upquote if available, for straight quotes in verbatim environments
\IfFileExists{upquote.sty}{\usepackage{upquote}}{}
\IfFileExists{microtype.sty}{% use microtype if available
  \usepackage[]{microtype}
  \UseMicrotypeSet[protrusion]{basicmath} % disable protrusion for tt fonts
}{}
\makeatletter
\@ifundefined{KOMAClassName}{% if non-KOMA class
  \IfFileExists{parskip.sty}{%
    \usepackage{parskip}
  }{% else
    \setlength{\parindent}{0pt}
    \setlength{\parskip}{6pt plus 2pt minus 1pt}}
}{% if KOMA class
  \KOMAoptions{parskip=half}}
\makeatother
\usepackage{xcolor}
\usepackage[margin=1in]{geometry}
\usepackage{color}
\usepackage{fancyvrb}
\newcommand{\VerbBar}{|}
\newcommand{\VERB}{\Verb[commandchars=\\\{\}]}
\DefineVerbatimEnvironment{Highlighting}{Verbatim}{commandchars=\\\{\}}
% Add ',fontsize=\small' for more characters per line
\newenvironment{Shaded}{}{}
\newcommand{\AlertTok}[1]{\textcolor[rgb]{1.00,0.00,0.00}{\textbf{#1}}}
\newcommand{\AnnotationTok}[1]{\textcolor[rgb]{0.38,0.63,0.69}{\textbf{\textit{#1}}}}
\newcommand{\AttributeTok}[1]{\textcolor[rgb]{0.49,0.56,0.16}{#1}}
\newcommand{\BaseNTok}[1]{\textcolor[rgb]{0.25,0.63,0.44}{#1}}
\newcommand{\BuiltInTok}[1]{\textcolor[rgb]{0.00,0.50,0.00}{#1}}
\newcommand{\CharTok}[1]{\textcolor[rgb]{0.25,0.44,0.63}{#1}}
\newcommand{\CommentTok}[1]{\textcolor[rgb]{0.38,0.63,0.69}{\textit{#1}}}
\newcommand{\CommentVarTok}[1]{\textcolor[rgb]{0.38,0.63,0.69}{\textbf{\textit{#1}}}}
\newcommand{\ConstantTok}[1]{\textcolor[rgb]{0.53,0.00,0.00}{#1}}
\newcommand{\ControlFlowTok}[1]{\textcolor[rgb]{0.00,0.44,0.13}{\textbf{#1}}}
\newcommand{\DataTypeTok}[1]{\textcolor[rgb]{0.56,0.13,0.00}{#1}}
\newcommand{\DecValTok}[1]{\textcolor[rgb]{0.25,0.63,0.44}{#1}}
\newcommand{\DocumentationTok}[1]{\textcolor[rgb]{0.73,0.13,0.13}{\textit{#1}}}
\newcommand{\ErrorTok}[1]{\textcolor[rgb]{1.00,0.00,0.00}{\textbf{#1}}}
\newcommand{\ExtensionTok}[1]{#1}
\newcommand{\FloatTok}[1]{\textcolor[rgb]{0.25,0.63,0.44}{#1}}
\newcommand{\FunctionTok}[1]{\textcolor[rgb]{0.02,0.16,0.49}{#1}}
\newcommand{\ImportTok}[1]{\textcolor[rgb]{0.00,0.50,0.00}{\textbf{#1}}}
\newcommand{\InformationTok}[1]{\textcolor[rgb]{0.38,0.63,0.69}{\textbf{\textit{#1}}}}
\newcommand{\KeywordTok}[1]{\textcolor[rgb]{0.00,0.44,0.13}{\textbf{#1}}}
\newcommand{\NormalTok}[1]{#1}
\newcommand{\OperatorTok}[1]{\textcolor[rgb]{0.40,0.40,0.40}{#1}}
\newcommand{\OtherTok}[1]{\textcolor[rgb]{0.00,0.44,0.13}{#1}}
\newcommand{\PreprocessorTok}[1]{\textcolor[rgb]{0.74,0.48,0.00}{#1}}
\newcommand{\RegionMarkerTok}[1]{#1}
\newcommand{\SpecialCharTok}[1]{\textcolor[rgb]{0.25,0.44,0.63}{#1}}
\newcommand{\SpecialStringTok}[1]{\textcolor[rgb]{0.73,0.40,0.53}{#1}}
\newcommand{\StringTok}[1]{\textcolor[rgb]{0.25,0.44,0.63}{#1}}
\newcommand{\VariableTok}[1]{\textcolor[rgb]{0.10,0.09,0.49}{#1}}
\newcommand{\VerbatimStringTok}[1]{\textcolor[rgb]{0.25,0.44,0.63}{#1}}
\newcommand{\WarningTok}[1]{\textcolor[rgb]{0.38,0.63,0.69}{\textbf{\textit{#1}}}}
\setlength{\emergencystretch}{3em} % prevent overfull lines
\providecommand{\tightlist}{%
  \setlength{\itemsep}{0pt}\setlength{\parskip}{0pt}}
\setcounter{secnumdepth}{-\maxdimen} % remove section numbering
\setlength{\emergencystretch}{3em}
\ifLuaTeX
  \usepackage{selnolig}  % disable illegal ligatures
\fi
\usepackage{bookmark}
\IfFileExists{xurl.sty}{\usepackage{xurl}}{} % add URL line breaks if available
\urlstyle{same}
\hypersetup{
  colorlinks=true,
  linkcolor={blue},
  filecolor={Maroon},
  citecolor={Blue},
  urlcolor={blue},
  pdfcreator={LaTeX via pandoc}}

\author{}
\date{}

\begin{document}

\section{DDTA research work plan after documentation-layer
closure}\label{ddta-research-work-plan-after-documentation-layer-closure}

\textbf{WORK PLAN - REVISION 8}

\textbf{Status:} active execution plan after BA1 minimal
identity-ontology closure.

\textbf{Supersedes:} Revision 7 only for forward execution state; R1-R7
remain historical research records.

\subsection{1. Fixed prior state}\label{fixed-prior-state}

\begin{itemize}
\tightlist
\item
  Chapters 2-4: CLOSED / FINAL for current thesis scope.
\item
  Documentation layer: CLOSED.
\item
  BA0-R systems-modeling prior-art research: CLOSED.
\item
  BA0 responsibility/non-goals: CLOSED by BA0-T3.
\item
  BA1 minimal BAE identity ontology: CLOSED by BA1-T3.
\item
  ThreatForge is a case study/tooling instrument, not DDTA semantic
  authority.
\end{itemize}

\subsection{2. Closed BA1 identity
ontology}\label{closed-ba1-identity-ontology}

BA1 accepts exactly:

\begin{Shaded}
\begin{Highlighting}[]
\NormalTok{BAReferent}
\NormalTok{BAProposition}
\end{Highlighting}
\end{Shaded}

\texttt{BAReferent} is the identity family for independently reusable
methodology-neutral shared project meaning.

\texttt{BAProposition} is the identity family for independently
reviewable methodology-neutral analytical assertions about BAReferents.

\texttt{BAE} is an umbrella term, not an additional required metaclass.

A BAReferent may denote structural, human/external,
informational/resource, behavioral/event, contractual, storage,
boundary, state/context or other project meaning. Current evidence does
not require dedicated first-class roots for those categories.

\subsection{3. Closed specialization
rule}\label{closed-specialization-rule}

A new dedicated BAE type/family is justified only if evidence shows
both:

\begin{enumerate}
\def\labelenumi{\arabic{enumi}.}
\tightlist
\item
  independent semantic identity across propositions/projections/change;
  and
\item
  reusable subtype-specific invariants that cannot be represented
  honestly as classifications, roles, values or propositions over the
  closed families.
\end{enumerate}

This prevents future BA phases from turning every useful semantic
category into a metaclass by convenience.

\subsection{4. Phase sequence after BA1}\label{phase-sequence-after-ba1}

Do not skip phases:

\begin{Shaded}
\begin{Highlighting}[]
\NormalTok{BA2 {-} Relations and canonical action vocabulary}
\NormalTok{BA3 {-} Document{-}to{-}BAE derivation, provenance and authority}
\NormalTok{BA4 {-} Multi{-}level human and method projections}
\NormalTok{BA5 {-} Controlled vocabulary and optional authoring/extraction assistance}
\NormalTok{BA6 {-} Base Analysis closure and regression}
\end{Highlighting}
\end{Shaded}

Formal analysis-envelope and threat-method work remains after BA6.

\subsection{5. BA2 objective}\label{ba2-objective}

BA2 decides the smallest stable methodology-neutral
proposition/relation/action structure over the closed BA1 identity
ontology.

BA2 must preserve three separations:

\begin{Shaded}
\begin{Highlighting}[]
\NormalTok{semantic identity family}
\NormalTok{    !=}
\NormalTok{semantic relation / participation role}
\NormalTok{    !=}
\NormalTok{lexical label / authoring synonym}
\end{Highlighting}
\end{Shaded}

It must also keep method-owned threat semantics outside Base Analysis.

BA2 may define how reusable semantic classifications are represented,
but it may not silently introduce new BA1 first-class identity families.
A genuine third-family requirement reopens BA1.

\subsection{6. BA2-T1 - proposition shape and participation-role lower
bound}\label{ba2-t1---proposition-shape-and-participation-role-lower-bound}

\textbf{NEXT / NOT EXECUTED BY THIS PLAN REVISION.}

BA2-T1 asks:

\begin{quote}
What is the minimum methodology-neutral structural shape required for
one BAProposition to express reusable project facts over BAReferents
without freezing an exhaustive predicate vocabulary or importing a
domain/method taxonomy?
\end{quote}

\subsubsection{Pressure questions}\label{pressure-questions}

BA2-T1 must test at least:

\begin{enumerate}
\def\labelenumi{\arabic{enumi}.}
\tightlist
\item
  whether binary subject-predicate-object is sufficient or n-ary
  participation is required;
\item
  how one proposition can represent roles such as source, destination,
  information, correlation or responsibility without turning each role
  into a BA1 type;
\item
  whether action/behavior identity belongs to BAReferent while action
  semantics belong to proposition structure;
\item
  how qualifiers such as condition, state, failure, applicability or
  realization attach without creating assertion ambiguity;
\item
  how the same project meaning can participate in different semantic
  roles across propositions;
\item
  whether method-neutral classification must be explicit for mechanical
  projection or can be derived from proposition roles without raw-prose
  reconstruction.
\end{enumerate}

\subsubsection{Guardrails}\label{guardrails}

BA2-T1 must not:

\begin{itemize}
\tightlist
\item
  invent an exhaustive verb list;
\item
  import STRIDE/DFD categories;
\item
  define BA3 provenance fields;
\item
  define BA4 views;
\item
  create ThreatForge implementation classes;
\item
  reopen BA1 unless a concrete identity counterexample is found.
\end{itemize}

\subsubsection{Exit condition}\label{exit-condition}

Produce a falsifiable lower-bound proposition structure with explicit
alternatives rejected/deferred. Do not yet claim a complete canonical
relation/action vocabulary unless the evidence genuinely forces it.

\subsection{7. Later BA2 work - provisional sequence
only}\label{later-ba2-work---provisional-sequence-only}

Only after BA2-T1 is reviewed should later BA2 microsteps pressure-test
and close the actual relation/action vocabulary.

Likely concerns include:

\begin{itemize}
\tightlist
\item
  stable relation/action semantic identities;
\item
  role cardinalities and n-ary propositions;
\item
  dependency, responsibility, correlation, realization and constraint
  semantics;
\item
  classification versus relation boundaries;
\item
  lexical canonical labels versus authoring synonyms.
\end{itemize}

This section is planning only, not semantic acceptance.

\subsection{8. BA3 - derivation, provenance and
authority}\label{ba3---derivation-provenance-and-authority}

After BA2 closure, BA3 must close:

\begin{itemize}
\tightlist
\item
  document-to-BAReferent/BAProposition derivation rules;
\item
  source locators;
\item
  grounded, derived and diagnostic/unresolved materialization;
\item
  identity/equivalence/lifecycle rules;
\item
  acceptance/rejection/stale semantics;
\item
  source-to-analysis and analysis-to-source traceability.
\end{itemize}

Governed documentation remains project authority.

\subsection{9. BA4 - projections}\label{ba4---projections}

Define derived human and method projections from the same accepted BA
semantics. No view becomes source authority.

\subsection{10. BA5 - vocabulary and optional
assistance}\label{ba5---vocabulary-and-optional-assistance}

Define controlled vocabulary and optional assistance only after the
semantic contracts exist. Lexical normalization, NLP/LLM proposals and
authoring suggestions remain support mechanisms and cannot establish
project semantics silently.

\subsection{11. BA6 - Base Analysis closure and
regression}\label{ba6---base-analysis-closure-and-regression}

Regress the implemented Base Analysis design against closed corpora and
at least one structurally different holdout. BA6, not BA1, closes the
complete Base Analysis design.

\subsection{12. Analysis envelope remains
downstream}\label{analysis-envelope-remains-downstream}

Only after Base Analysis closure:

\begin{itemize}
\tightlist
\item
  A1 - AnalysisRecord / execution envelope;
\item
  A2 - common reviewed Finding boundary;
\item
  A3 - change/provenance integration;
\item
  formal methodology overlays and STRIDE/STRIDE-AI evaluations.
\end{itemize}

The BA0-T2 STRIDE probe remains only a pre-overlay consumer probe.

\subsection{13. Next authorized
microstep}\label{next-authorized-microstep}

After the BA1-T3 closure package is committed and remotely verified,
execute only:

\begin{quote}
\textbf{BA2-T1 - minimal BAProposition structural shape and
participation-role lower-bound derivation.}
\end{quote}

Do not start formal STRIDE overlay design, Common Finding schema or
implementation work.

\end{document}
